\documentclass[11pt]{article}
\usepackage[margin=1in]{geometry}
\usepackage{hyperref}
\usepackage{listings}
\usepackage{xcolor}
\usepackage{amsmath}

\lstset{
  basicstyle=\ttfamily\small,
  breaklines=true,
  columns=fullflexible,
  frame=single,
  backgroundcolor=\color{gray!10},
  showstringspaces=false
}

\title{URI OPeNDAP Test Case: Working Log (Prompt \#3)}
\author{Compiled for P. Cornillon}
\date{2026-05-16}

\begin{document}
\maketitle

\section*{Purpose}

Working log for Prompt~\#3 of the URI OPeNDAP test case (the
\texttt{gradients\_by\_period/} analog of Prompt~\#2's
\texttt{SST\_Orbits/} audit). Findings are in
\texttt{docs/uri\_test\_case\_3.tex}.

\section{Prompt \#3 (2026-05-16 20:26 UTC)}

\subsection{Instruction}

\emph{``Repeat the above for Readable Branch \#2:
gradients\_by\_period/. Since this archive has no semantic metadata,
you will have to infer whatever you can. Group the semantic
metadata by inferred values for which you are very confident, those
for which you are not sure but have guessed at entries, and those
for which you do not feel comfortable guessing. Also group by
essential and `would be nice'.''}

\subsection{Starting point from Prompts~\#1--2}

\begin{itemize}
  \item 270 files, one per calendar month from \texttt{01\_2003}
    to \texttt{12\_2024}; filename pattern
    \texttt{augmented\_monthly\_stats\_MM\_YYYY.nc}.
  \item DMR: 2 dimensions (\texttt{lon=360}, \texttt{lat=180}),
    34 \texttt{Float64} variables of shape (lon, lat); zero
    global attributes; zero variable attributes.
  \item Variable naming grammar is regular:
    \texttt{day\_/night\_} prefix; one of six quantity families
    (SST, eastward/northward gradient, magnitude gradient, and
    swath-relative grad\_as/grad\_at/grad\_mag\_per\_km); for sums,
    a paired \texttt{\_squared}; three pixel-count flavors
    (\texttt{pixel\_count}, \texttt{as\_pixel\_count},
    \texttt{at\_pixel\_count}).
  \item The companion branch \texttt{SST\_Orbits/} carries
    \texttt{SST\_In} in $^\circ$C and
    \texttt{eastward\_gradient/northward\_gradient} in $^\circ$C/km;
    that is the source for the unit assignments here.
\end{itemize}

\subsection{New investigations}

Three additional probes were run for Prompt~\#3.

\paragraph{Schema constancy across the period.}
\texttt{scripts/gradients\_dmr\_inventory.py} fetches the DMR for
\texttt{01\_2003}, \texttt{07\_2020}, and \texttt{12\_2024} and
confirms the schema is identical: 2 dimensions
(\texttt{lon=360}, \texttt{lat=180}), 34 Float64 variables of
shape (lon, lat), no attributes anywhere. Conclusion: assumptions
made from a single sample file generalize across all 270 files.

\paragraph{Lat / lon convention by geography.}
\texttt{scripts/gradients\_geom\_resolve.py} tests the four
combinations of \{lat axis runs S--N or N--S\} $\times$
\{lon axis is $-180\to+180$ or $0\to360$\} by looking up six
known-land points (Sahara, Australia interior, Amazon basin,
Central US, Greenland, Tibet) and six known-ocean points (Pacific
center, Atlantic center, Indian Ocean, North Pacific, Persian Gulf,
Western Pacific Warm Pool) and asking whether \texttt{day\_pixel\_count}
is zero where it should be (land) and non-zero where it should be
(ocean). The result is unambiguous: only the combination
\textbf{lat\_idx 0 = $-89.5^\circ$} (S--N), \textbf{lon\_idx 0 =
$-179.5^\circ$} (cell-centered $-180\to+180$) classifies all 12
points correctly. All other conventions miss multiple points
(e.g.\ N--S lat would put Australia at an ocean-only index).

The implied cell-center coordinate vectors are
\begin{lstlisting}
lat[k] = -89.5 + k       (k = 0 .. 179)
lon[k] = -179.5 + k      (k = 0 .. 359)
\end{lstlisting}

\paragraph{Sign and divisor conventions on gradient variables.}
\texttt{scripts/gradients\_signs\_test.py} pulls the day-side
gradient sums and squared sums for \texttt{07\_2020} and reports:
\begin{itemize}
  \item \texttt{eastward\_gradient}, \texttt{northward\_gradient},
    \texttt{grad\_as\_per\_km}, \texttt{grad\_at\_per\_km} are
    all signed (both positive and negative sums observed).
  \item \texttt{magnitude\_gradient} and \texttt{grad\_mag\_per\_km}
    are strictly non-negative (min $= 0$).
  \item The squared sums are all non-negative, as required.
  \item \texttt{sum\_magnitude\_gradient} per cell is much larger
    than $\sqrt{S_e^2 + S_n^2}$ where $S_e =
    \mathtt{sum\_eastward\_gradient}$ and $S_n =
    \mathtt{sum\_northward\_gradient}$ in the same cell; the ratio
    mean is $\approx 12$. That tells us
    \texttt{sum\_magnitude\_gradient} is the sum of per-pixel
    magnitudes $\sum |\nabla T|$ rather than the magnitude of the
    vector sum (which would yield ratio $= 1$). Standard
    bookkeeping.
  \item Pixel-count comparisons: \texttt{day\_pixel\_count} has
    mean 10\,030 per cell, \texttt{day\_as\_pixel\_count} has
    mean 8\,110 ($\approx 0.81\,N_{\rm pix}$), and
    \texttt{day\_at\_pixel\_count} has mean 4\,090
    ($\approx 0.41\,N_{\rm pix}$). The three counts differ:
    along-scan gradient validity is $\sim$2$\times$ the along-track
    gradient validity. The latter ratio is consistent with MODIS
    Aqua's 10-line scan structure, where every tenth along-track
    derivative crosses a scan boundary and is masked.
  \item Per-cell mean check: $\mathtt{sum\_magnitude\_gradient}
    / \mathtt{pixel\_count}$ gives $\sim 0.046\,^\circ$C/km
    averaged globally; $\mathtt{sum\_grad\_mag\_per\_km}
    / \mathtt{at\_pixel\_count}$ gives $\sim 0.044\,^\circ$C/km;
    $\mathtt{sum\_grad\_mag\_per\_km} / \mathtt{as\_pixel\_count}$
    gives $\sim 0.022\,^\circ$C/km. The first two match closely,
    indicating that \texttt{sum\_grad\_mag\_per\_km} should be
    divided by \texttt{at\_pixel\_count} (the more restrictive
    count), not \texttt{as\_pixel\_count}.
\end{itemize}

Combined, this fixes the per-variable divisor map:
\begin{itemize}
  \item \texttt{*\_sum\_SST}, \texttt{*\_sum\_SST\_squared} :
    divide by \texttt{*\_pixel\_count}.
  \item \texttt{*\_sum\_eastward\_gradient}, etc.,
    \texttt{*\_sum\_magnitude\_gradient}, etc. : divide by
    \texttt{*\_pixel\_count}.
  \item \texttt{*\_sum\_grad\_as\_per\_km}, etc.: divide by
    \texttt{*\_as\_pixel\_count}.
  \item \texttt{*\_sum\_grad\_at\_per\_km}, etc.: divide by
    \texttt{*\_at\_pixel\_count}.
  \item \texttt{*\_sum\_grad\_mag\_per\_km}, etc.: divide by
    \texttt{*\_at\_pixel\_count} (inferred from value match).
\end{itemize}

\subsection{What I am NOT able to pin down}

\begin{itemize}
  \item The exact \emph{algorithm} used in the upstream orbit
    files to compute gradients (forward difference, centered
    difference, Sobel-style operator, distance-weighted between
    neighboring pixels, etc.). The orbit-file standard\_names
    just say ``\texttt{eastward\_temperature\_gradient}'' and
    ``\texttt{northward\_temperature\_gradient}'' without
    specifying.
  \item Whether \texttt{magnitude\_gradient} and
    \texttt{grad\_mag\_per\_km} are computed pixel-by-pixel from
    the geographic and swath-relative components respectively,
    or whether one is just stored twice with different counts.
    At a pixel, $|\nabla T|$ is invariant under rotation, so the
    two should be numerically equal pixel-by-pixel; the file
    distinguishes them, so presumably they differ in
    \emph{which pixels qualify as valid}. That is consistent with
    the count comparison above.
  \item Whether `\texttt{as}'/`\texttt{at}' really stand for
    along-scan / along-track. The values are signed and have
    almost identical standard deviations to
    \texttt{eastward}/\texttt{northward}, so at low latitude they
    are nearly identical to the geographic components. The
    name strongly implies along-scan / along-track but I have
    not proved it.
  \item Provenance attributes (\texttt{date\_created},
    \texttt{history}, \texttt{processing\_software}, exact UUID
    of the generating run, etc.) cannot be inferred from the
    file or the directory listing.
\end{itemize}

\subsection{Plan for the findings document}

The findings document organizes inferred values along two axes
as the user requested:
\begin{itemize}
  \item \textbf{Importance.} \emph{Essential} = required by COARDS
    or by any reasonable reading of the data (units, fill,
    coordinate variables, basic globals). \emph{Nice-to-have} =
    CF/ACDD attributes that aid interpretation but are not
    required.
  \item \textbf{Confidence.} \emph{Very confident} = directly
    derivable from data, naming convention, or the companion
    branch's metadata. \emph{Guessed} = plausible best guess that
    a reader could not verify without contacting the producer.
    \emph{Will not guess} = the gap exists but I do not have
    enough evidence to propose a value.
\end{itemize}

The cross of those gives a $3 \times 2$ table around which the
audit is structured.

\subsection{Commands}

\begin{lstlisting}
python3 scripts/gradients_dmr_inventory.py
python3 scripts/gradients_geom_resolve.py
python3 scripts/gradients_signs_test.py
\end{lstlisting}

\subsection{Session timing}

Started 2026-05-16 20:26\,UTC. Closed log 2026-05-16
$\sim$20:55\,UTC.

\end{document}
